\chapter{Referencial Te\'orico}
\label{pre:referencial}

A evas\~ao escolar e o desempenho acad\^emico abaixo do esperado em institui\c{c}\~oes de ensino prejudicam n\~ao apenas o progresso educacional dos alunos, mas tamb\'em a efic\'acia dos sistemas de ensino. Al\'em disso, a dificuldade em identificar e intervir precocemente em problemas acad\^emicos limita a capacidade das institui\c{c}\~oes de proporcionar uma educa\c{c}\~ao de qualidade.

O objetivo deste trabalho \'e explorar o uso de redes neurais de grafos na an\'alise de dados educacionais como uma abordagem inovadora para enfrentar o problema da evas\~ao escolar e melhorar o desempenho dos estudantes. As redes neurais de grafos emergiram como uma alternativa promissora devido \`a sua capacidade de modelar rela\c{c}\~oes complexas entre entidades interconectadas.

\section{Redes Neurais de Grafos}
\label{pre:rng}

As redes neurais de grafos s\~ao uma extens\~ao das redes neurais convencionais voltada \`a modelagem de dados interconectados. Essa abordagem representa os dados como um grafo em que os n\'os representam entidades e as arestas capturam as rela\c{c}\~oes entre elas. Dessa forma, essas redes conseguem aprender representa\c{c}\~oes ricas e hier\'arquicas dos dados.

Casalegno \cite{casalegno2021} destaca que grafos s\~ao estruturas altamente vers\'ateis para aprendizado de m\'aquina, mas que seu tratamento n\~ao \'e trivial por envolver dados n\~ao estruturados e n\~ao euclidianos. Nesse cen\'ario, as redes convolucionais de grafos (\sigla{GCN}{Graph Convolutional Networks}) ampliam o conceito das redes neurais convolucionais para dados estruturados em forma de grafo, permitindo tarefas como classifica\c{c}\~ao de grafos, classifica\c{c}\~ao de n\'os e previs\~ao de liga\c{c}\~oes.

\section{An\'alise de Dados Educacionais}
\label{pre:analise}

A an\'alise de dados educacionais envolve a coleta e o processamento de informa\c{c}\~oes relacionadas ao desempenho acad\^emico dos alunos, intera\c{c}\~oes em sala de aula e outros fatores relevantes. Seu objetivo \'e fornecer \textit{insights} que possam melhorar a qualidade do ensino, identificar padr\~oes de desempenho e prever evas\~ao escolar.

Dekker, Pechenizkiy e Vleeshouwers \cite{dekker2009} demonstraram, em investiga\c{c}\~ao de minera\c{c}\~ao de dados educacionais, que classificadores relativamente simples e de f\'acil interpreta\c{c}\~ao, como \'arvores de decis\~ao, podem produzir resultados valiosos para antecipar evas\~ao de estudantes, alcan\c{c}ando precis\~ao entre 75\% e 80\%.

\section{Redes Neurais de Grafos na Educa\c{c}\~ao}
\label{pre:educacao}

O uso de redes neurais de grafos na an\'alise de dados educacionais ganhou destaque por sua capacidade de lidar com a complexidade inerente aos sistemas educacionais. Essas redes podem capturar relacionamentos multifacetados entre alunos, cursos, professores e outras entidades envolvidas no processo educacional. Pesquisas anteriores apontam aplica\c{c}\~oes como:

\begin{enumerate}
    \item detec\c{c}\~ao de padr\~oes de comportamento do aluno;
    \item previs\~ao de abandono escolar;
    \item recomenda\c{c}\~ao de cursos e recursos;
    \item avalia\c{c}\~ao de efic\'acia do ensino.
\end{enumerate}

\section{Comparando com m\'etodos tradicionais}
\label{pre:comparacao}

Embora os m\'etodos tradicionais de an\'alise de dados educacionais sejam amplamente utilizados, eles geralmente apresentam limita\c{c}\~oes na captura de rela\c{c}\~oes complexas. Bhatti et al. \cite{bhatti2023} apontam que redes neurais de grafos, adaptadas para dados n\~ao euclidianos, t\^em sido amplamente usadas em diversas \'areas justamente por sua capacidade de lidar com estruturas em grafo. Essa vantagem refor\c{c}a o interesse em investigar sua aplica\c{c}\~ao em cen\'arios educacionais.

Essa revis\~ao estabelece a base do estudo proposto e sustenta a investiga\c{c}\~ao sobre como redes neurais de grafos poderiam ser empregadas para melhorar a compreens\~ao de padr\~oes acad\^emicos e prevenir a evas\~ao escolar.
